%%%%%%%%%%%%%%%%
%%% gen 12 %%%
%%%%%%%%%%%%%%%%

\Chapter{12}

\Verse{1}
{
	\hJ{וַיֹּאמֶר יְהֹוָה אֶל אַבְרָם לֶךְ לְךָ מֵאַרְצְךָ וּמִמּוֹלַדְתְּךָ וּמִבֵּית אָבִיךָ אֶל הָאָרֶץ אֲשֶׁר אַרְאֶךָּ׃}
}
{
	\eJ{
		And YHWH said to Abram, “Go from your land and from your
		birthplace and from your father’s house to the land that
		I’ll show you.
	}
}
{
	Ok so YHWH tells Abe to leave home and go somewhere.

	He hasn't decided where he's sending him yet.

	But it's time to leave home.

	Which he already did in Gen 11:31.\footnote{Give YHWH a break, he can't know everything.}

%	\footnote{
%		Go. Hebrew lek lk. Much has been made of the second word in this phrase,
%		lk, which means “for you.” No translation quite captures the sense of
%		the Hebrew (“Go you,” “Get you,” “Go for yourself”). The second element,
%		lk (or lkem), occasionally follows verbs in the Tanak. For example, in
%		Lev 26:1 we have l’ ta‘ lkem (you shall not make for yourselves; you
%		shall make you no …), and in Num 14:25 we have s‘ lkem (travel you;
%		travel for yourselves). I believe it is better to use no English term
%		than to use any of the possible equivalents, all of which are clumsy
%		English and do not convey the Hebrew. Footnote 12:1. your land, your
%		birthplace. But he is in Haran. He has already left Ur of the Chaldeans,
%		which was his land and birthplace. Rashi tries to solve the problem by
%		saying that this means “Go even farther from your land and birthplace.”
%		Ibn Ezra tries to solve it by saying that it means “God had said [i.e.,
%		earlier, back in Ur], ‘Go from your land …’” But these stretch the
%		meaning of the Hebrew beyond what is possible. And besides, Abraham
%		later sends his servant to get a wife for Isaac, saying, “Go to my land
%		and my birthplace” (24:4), and the servant goes to Haran (“to Aram
%		Naharaim, to the city of Nahor,” 24:10; which is identified as Haran,
%		27:43). On these grounds Ramban rejects Rashi’s and Ibn Ezra’s
%		solutions. Ramban’s own solution, though, falls short as well, requiring
%		that Abraham’s land and birthplace are Haran, not Ur. This is further
%		complicated by the fact that the Abrahamic covenant begins with God
%		saying to Abraham, “I brought you out of Ur …” (15:7). The truth is that
%		this is a case in which the contradiction is a result of the fact that
%		the Torah was composed from several sources. In this case, one of its
%		sources had Abraham coming from Ur, and another had him coming from
%		Haran. The aim of this commentary is to deal with the Torah in its
%		final, combined form and not to analyze its sources. (I have done that
%		elsewhere: in Who Wrote the Bible? and The Hidden Book in the Bible.)
%		Usually, the combination of sources has produced a richer, more complex
%		story in the Torah. But this passage is one of a few instances in which
%		the combination has produced an irreconcilable contradiction. In such
%		cases it is better to acknowledge the problem than to stretch to make
%		forced interpretations. Footnote 12:1. your land, your birthplace, and
%		your father’s house. But if he has left his land, then of course he has
%		left his birthplace and his father’s house. This is geographically
%		backwards. Therefore, the point of this order is not geographical. It is
%		emotional. The three steps are arranged in ascending order of difficulty
%		for Abraham. It is hard to leave one’s land, harder if it is where one
%		was born, and harder still to leave one’s family. And where is he to go?
%		To “the land that I’ll show you.” That is, he must leave his homeland
%		without knowing for what he is giving it up. The wording seems designed
%		to make it hard for Abraham. This pattern of testing will be repeated
%		when he is commanded to sacrifice Isaac, thus: “Take your son, your only
%		one, whom you love, Isaac.” And where is he to sacrifice him? “On one of
%		the mountains that I’ll say to you.” (See the comment on 22:2.)
%	}
}

\Verse{2}
{
	\hJ{וְאֶעֶשְׂךָ לְגוֹי גָּדוֹל וַאֲבָרֶכְךָ וַאֲגַדְּלָה שְׁמֶךָ וֶהְיֵה בְּרָכָה׃}
}
{
	\eJ{
		And I’ll make you into a big nation, and I’ll bless you
		and make your name great. And be a blessing!
	}
}
{
	He promises to make him into a country. And to make his name great.\fA{His current name is apparently not great, which may be related to why YHWH changes it later.}
}

\Verse{3}
{
	\hJ{וַאֲבָרְכָה מְבָרְכֶיךָ וּמְקַלֶּלְךָ אָאֹר וְנִבְרְכוּ בְךָ כֹּל מִשְׁפְּחֹת הָאֲדָמָה׃}
}
{
	\eJ{
		And I’ll bless those who bless you, and those who affront
		you I’ll curse. And all the families of the earth will be
		blessed through you.”
	}
}
{
	He promises a lot.
	
	Your friends will be my friends.
	
	Your enemies will be my enemies.

	And all the families of the earth from all the other languages and countries, they're gonna be blessed too.%
	\footnote{They exist now, because of Gen. 11, when YHWH invented languages and countries out of compassion.}

	Because of Abe.

%	\footnote{
%		all the families of the earth will be blessed through you. The Tanak is
%		so focused on the people of Israel that one can underestimate the
%		overwhelming significance of its opening section. The first eleven
%		chapters (Parashat Bereshit and Parashat Noah) are about the
%		relationship between God and the entire human community. That
%		relationship does not go well, and after ten generations the deity
%		decides to destroy the mass and start over with a single virtuous man’s
%		family. But it turns out that choosing a virtuous individual does not
%		guarantee that this individual’s descendants will be virtuous as well.
%		Another ten generations pass, and humans in general are not a
%		planet-full of Noahs. So once again the focus narrows to a single
%		virtuous person, Abraham. We must keep in mind what has happened up to
%		this point when we read this, or else we will lose the significance of
%		what is happening here. Wiping out everyone but a virtuous person did
%		not work. So God leaves the species alive but chooses an individual who
%		will produce a family that will ultimately bring blessing to all the
%		families of the earth. To make sure we get it, it is the final point of
%		God’s first revelation to Abraham: Go to the land I’ll show you. I’ll
%		make you a big nation. I’ll bless you. “And all the families of the
%		earth will be blessed through you” (Gen 12:3). To make sure we do not
%		forget it, it is repeated four times—and always in crucial moments of
%		revelation: during the appearance of the three visitors to Abraham (Gen
%		18:17–18); in the blessing following the near-sacrifice of Isaac
%		(22:16–18); in God’s first appearance to Isaac (26:2–4); and in Jacob’s
%		first encounter with God in the dream of the ladder at Beth-El
%		(28:10–14). It is important. In some way, at some time, the result of
%		the divine choice of Abraham is supposed to be some good for all
%		humankind. We are never told what this good is supposed to be. Is it
%		that Abraham’s descendants are supposed to bring blessing by being “a
%		light to the nations”—setting an example, showing how a community can
%		live: caring for one another, not cheating one another, not enslaving
%		one another, not lending to each other for profit, and so on and on? Or
%		is it that they are to bring blessing by doing things that benefit the
%		species: inventions, cures, literature, music, learning? It does not
%		say. But at minimum it must mean that the people of Israel do not live
%		alone or apart. Their destiny—our destiny—whatever it is, must be bound
%		up in the destiny of all humankind. This adds a dimension, an additional
%		layer of significance, to every story that will follow in the Torah.
%		When Abraham travels to Canaan, we might imagine it from a God’s-eye
%		view above the earth: the tiny movement of a man and his family along
%		the globe is a first step in a process that is to bring benefit for all
%		the earth. When Abraham travels to Egypt, this is the first in a series
%		of encounters that he and his descendants will have with the peoples of
%		the world. When he and Lot are forced to part, this is a step in
%		distinguishing his destiny from that of the families of his brothers
%		Haran and Nahor, though he and his descendants will continue to interact
%		with them. When he joins in a battle among kings in order to rescue Lot
%		(Genesis 14), he is being drawn into world events, and it is another
%		first step, an anticipation of all the times that his descendants, the
%		people of Israel, will be drawn into contact and interaction with
%		nations. When Abraham covenants with God (Genesis 15 and 17) the
%		ceremony includes specific references to Ur of the Chaldees, the land
%		from which Abraham has come; to Egypt, the land where his descendants
%		will be enslaved; and to ten peoples of the land that Abraham is
%		promised. It also includes the announcement of his coming son Isaac, the
%		key to the fulfillment of the destiny to be a blessing to all the
%		families of the earth—as will be confirmed explicitly following the
%		Aqedah. And in between the two chapters concerning the Abrahamic
%		covenant comes the story of the birth of Abraham’s son Ishmael, whose
%		descendants, the Ishmaelites, will be among Israel’s related neighbors.
%		All of these episodes are conveying the formative stages of a
%		development that begins with that narrowing of attention to Abraham. The
%		interpretive point is that we must understand every section of the Torah
%		with awareness of what has preceded it and what will come after it. The
%		social and moral point is that Abraham’s descendants are not to live by
%		themselves or only for themselves. Whether dealing with non-Jews who
%		live in Israel or dealing with non-Jews who are their neighbors in the
%		countries in which they reside, the Jews are a community that connects
%		its birth with a prediction (?), a promise (?), an obligation (?), a
%		destiny (?) to be a blessing to them all. Footnote 12:3. will be blessed
%		through you. Rashi takes the plain meaning of this phrase to mean that
%		non-Israelites will bless their children with words such as “May you be
%		like Abraham.” His proof for this reading is the wording of Jacob’s
%		blessing of his grandsons Ephraim and Manasseh: “Israel will bless with
%		you, saying: ‘May God make you like Ephraim and Manasseh’” (Gen 48:20).
%		But that passage uses an active form of the verb (the Piel), whereas all
%		the occurrences of this phrase in terms of the nations being blessed use
%		forms that are passive or reflexive (Niphal and Hitpael), which is a
%		different matter. If we take the meaning as passive, then it is as I
%		have translated it here: “the earth’s families will be blessed through
%		you.” If we take it as reflexive, then it can mean “they will bless
%		themselves through you,” which is Rashi’s understanding; but a reflexive
%		can also mean “they will get blessing through you.” Interpreters are
%		split on this. I believe that context settles the question. The issue in
%		the story until now has been the course of relations between God and all
%		the families of the earth. Now God makes a special bond with Abraham’s
%		family and lets him know that this is for the eventual benefit of all
%		families. (See the preceding comment.)
%	}
}

\Verse{4}
{
	\hJ{וַיֵּלֶךְ אַבְרָם כַּאֲשֶׁר דִּבֶּר אֵלָיו יְהֹוָה וַיֵּלֶךְ אִתּוֹ לוֹט}
	\hP{וְאַבְרָם בֶּן חָמֵשׁ שָׁנִים וְשִׁבְעִים שָׁנָה בְּצֵאתוֹ מֵחָרָן׃}
}
{
	\eJ{And Abram went as YHWH had spoken to him, and Lot went with him.}
	\eP{And Abram was seventy-five years old when he went out from Haran.}
}
{
	So Abe leaves home.
	
	His nephew Lot comes too.
	
	Abe's 75 at this point.
%	\footnote{
%		Abram went as YHWH had spoken to him. Traditionally the emphasis has
%		been on faith as the mark of Abraham’s character. But the narrative
%		conveys far more that the mark of Abraham is obedience. Leave your land:
%		“And Abram went.” Sacrifice your son: “And Abram got up early.” The
%		tests that Abraham undergoes are not necessarily tests of faith but are
%		certainly tests of obedience. At this stage in relations between God and
%		humans, God appears to single out a human who will do what he is told.
%		This will change gradually in the Tanak, as humans grow and mature and
%		God cedes more responsibility for the world to them. But for now,
%		obedience is sought.
%	}
}

\Verse{5}
{
	\hP{וַיִּקַּח אַבְרָם אֶת שָׂרַי אִשְׁתּוֹ וְאֶת לוֹט בֶּן אָחִיו וְאֶת כּל רְכוּשָׁם אֲשֶׁר רָכָשׁוּ וְאֶת הַנֶּפֶשׁ אֲשֶׁר עָשׂוּ בְחָרָן וַיֵּצְאוּ לָלֶכֶת אַרְצָה כְּנַעַן וַיָּבֹאוּ אַרְצָה כְּנָעַן׃}
}
{
	\eP{
		And Abram took Sarai, his wife, and Lot, his brother’s
		son, and all their property that they had accumulated and
		the persons whom they had gotten in Haran; and they went
		out to go to the land of Canaan. And they came to the land
		of Canaan.
	}
}
{
	He takes wife and nephew and their stuff and the people they own.

	And they take their first steps on the long journey to Canaan, which we will now document in detail.

	They arrive.
}

\Verse{6}
{
	\hJ{וַיַּעֲבֹר אַבְרָם בָּאָרֶץ עַד מְקוֹם שְׁכֶם עַד אֵלוֹן מוֹרֶה וְהַכְּנַעֲנִי אָז בָּאָרֶץ׃}
}
{
	\eJ{
		And Abram passed through the land as far as the place of
		Shechem, as far as the oak of Moreh. And the Canaanite was
		in the land then.
	}
}
{
	Abe passes through the land.
	
	As far as the future capital of the North Kingdom of Israel.

	As far as that tree, you know the one.

	Oh and the Canaanite was in the land then.
}

\Verse{7}
{
	\hJ{וַיֵּרָא יְהֹוָה אֶל אַבְרָם וַיֹּאמֶר לְזַרְעֲךָ אֶתֵּן אֶת הָאָרֶץ הַזֹּאת וַיִּבֶן שָׁם מִזְבֵּחַ לַיהֹוָה הַנִּרְאֶה אֵלָיו׃}
}
{
	\eJ{
		And YHWH appeared to Abram and said, “I’ll give this land
		to your seed.” And he built an altar there to YHWH who had
		appeared to him.
	}
}
{
	YHWH's back.

	He appears to Abe and says ``Abe! Babe.
	I'll give this land to your c\eR{hildren, }um, baby.''

	Which is apparently convincing.
	
	He builds an altar to YHWH.
}

\Verse{8}
{
	\hJ{וַיַּעְתֵּק מִשָּׁם הָהָרָה מִקֶּדֶם לְבֵית אֵל וַיֵּט אהֳלֹה בֵּית אֵל מִיָּם וְהָעַי מִקֶּדֶם וַיִּבֶן שָׁם מִזְבֵּחַ לַיהֹוָה וַיִּקְרָא בְּשֵׁם יְהֹוָה׃}
}
{
	\eJ{
		And he moved on from there to the hill country east of
		Beth-El and pitched his tent—Beth-El to the west and Ai to
		the east—and he built an altar there to YHWH and invoked
		the name YHWH.
	}
}
{
	Ok so he leaves the altar of YHWH and he goes to House of God.

	Actually he's between House of God and Heap of Ruins.
	
	He builds another altar and says ``YHWH'' again.
%	\footnote{
%		Beth-El to the west and Ai to the east. At this site, Abram invokes the
%		name YHWH for the first time, and he builds an altar there. Later
%		Abraham returns to this site and invokes the divine name again (13:3).
%		No explanation of the significance of this place between the two
%		cities—Beth-El and Ai—is given in the Torah. But we learn in the book of
%		Joshua that, many generations later, God delivers Ai into the
%		Israelites’ hands by means of an ambush, and it takes place “between
%		Beth-El and Ai, to the west of Ai” (Josh 8:9,12; cf. 12:9). This appears
%		to be a case of “the merit of the fathers.” That is, the pious acts of
%		Abraham, Isaac, and Jacob have consequences for their descendants. In
%		this case, the Israelites are successful in the age of settlement of the
%		land thanks to what Abraham had done on the very spot centuries earlier.
%		There will be enough cases of this phenomenon to establish it as a major
%		theme and message of the Torah: that the good acts that we do have
%		implications for the lives of our children and many generations
%		thereafter: as an example, as a source of strength, and in subtle ways
%		that are impossible to calculate or foresee. This phenomenon of the
%		merit of the parents also acts as a preparation for and a balance to the
%		opposite phenomenon that will be developed in the Torah later: that our
%		sins are visited on the lives of our children and descendants as well.
%	}
}

\Verse{9}
{
	\hJ{וַיִּסַּע אַבְרָם הָלוֹךְ וְנָסוֹעַ הַנֶּגְבָּה׃}
}
{
	\eJ{And Abram traveled on, going to the Negeb.}
}
{
	He goes to the Dry Place.
}

\Verse{10}
{
	\hJ{וַיְהִי רָעָב בָּאָרֶץ וַיֵּרֶד אַבְרָם מִצְרַיְמָה לָגוּר שָׁם כִּי כָבֵד הָרָעָב בָּאָרֶץ׃}
}
{
	\eJ{
		And there was a famine in the land, and Abram went down to
		Egypt to reside there because the famine was heavy in the
		land.
	}
}
{
	There's no food, so he goes to Egypt, because there's no food.
}

\Verse{11}
{
	\hJ{וַיְהִי כַּאֲשֶׁר הִקְרִיב לָבוֹא מִצְרָיְמָה וַיֹּאמֶר אֶל שָׂרַי אִשְׁתּוֹ הִנֵּה נָא יָדַעְתִּי כִּי אִשָּׁה יְפַת מַרְאֶה אָתְּ׃}
}
{
	\eJ{
		And it was when he was close to coming to Egypt: and he
		said to Sarai, his wife, “Here, I know that you’re a
		beautiful woman.
	}
%	\footnote{12:11–13. Here, I know. Abraham’s words
%		to Sarah contain the first two occurrences of the Hebrew
%		word na’. It is sometimes rendered “now” in English and
%		sometimes rendered “please.” But neither is quite right.
%		It is an untranslatable particle that is a sign of polite
%		speech. Its most striking occurrence is in God’s words to
%		Abraham in the story of the binding of Isaac: “Take na’
%		your son … and offer him” (Gen 22:2). I do not translate
%		it with any English word, because it is closer to the
%		Hebrew to use no word than to use such words as “please”
%		or “now.”
%	}
}
{
	Ok so when he was close to coming
	
	To Egypt, he says to his wife,
	
	``Look, you're hot.''
}

\Verse{12}
{
	\hJ{וְהָיָה כִּי יִרְאוּ אֹתָךְ הַמִּצְרִים וְאָמְרוּ אִשְׁתּוֹ זֹאת וְהָרְגוּ אֹתִי וְאֹתָךְ יְחַיּוּ׃}
}
{
	\eJ{
		And it will be when the Egyptians see you that they’ll
		say, ‘This is his wife,’ and they’ll kill me and keep you
		alive.
	}
}
{
	``The Egyptians are gonna take one look at you and kill me and keep you alive for, well, reasons.''
}

\Verse{13}
{
	\hJ{אִמְרִי נָא אֲחֹתִי אָתְּ לְמַעַן יִיטַב לִי בַעֲבוּרֵךְ וְחָיְתָה נַפְשִׁי בִּגְלָלֵךְ׃}
}
{
	\eJ{
		Say you’re my sister so it will be good for me on your
		account and I’ll stay alive because of you.”
	}
}
{
	Say you're my sister.
	
	It'll be good for me.
	
	And sure it won't exactly be good for---
	
	Scratch that.
	
	I'll stay alive because of you.
%	\footnote{
%		Say you’re my sister … and I’ll stay alive. All commentators have
%		agonized over this. Abraham constructs a lie and puts Sarah in a
%		compromising position. This may mean that he thinks that he simply has
%		no real choice: it is either lie or die. Or it may convey that
%		Abraham—like other biblical heroes—is not perfect. We cannot know. In
%		either of these cases, he cannot be faulted for choosing to put Sarah in
%		a compromising position, because, in his understanding, Sarah would be
%		taken either way. The concern is that, without the lie, they would also
%		kill him. (In fact he turns out to be wrong. There is no evidence in the
%		Bible that the Egyptians or anyone else ever did what Abraham fears they
%		would do. Here and in two other cases in which Abraham and his son Isaac
%		claim that their wives are their sisters, the Egyptian and Philistine
%		kings send the couple away when they find out that she is his wife.)
%		Footnote 12:13. my sister. For a while it was common to cite tablets
%		from Nuzi as evidence that there was a Near Eastern practice of adopting
%		one’s wife as a sister. But the Nuzi case involved a single family, and
%		scholars recently have questioned whether it is justified to make so
%		much of it. Abraham’s actions here (and in Genesis 20; and Isaac’s
%		actions in Genesis 26) are treated as an unusual case. They require no
%		Near Eastern legal precedent.
%	}
}

\Verse{14}
{
	\hJ{וַיְהִי כְּבוֹא אַבְרָם מִצְרָיְמָה וַיִּרְאוּ הַמִּצְרִים אֶת הָאִשָּׁה כִּי יָפָה הִוא מְאֹד׃}
}
{
	\eJ{
		And it was as Abram came to Egypt, and the Egyptians saw
		the woman, that she was very beautiful.
	}
}
{
	So that happens.
}

\Verse{15}
{
	\hJ{וַיִּרְאוּ אֹתָהּ שָׂרֵי פַרְעֹה וַיְהַלְלוּ אֹתָהּ אֶל פַּרְעֹה וַתֻּקַּח הָאִשָּׁה בֵּית פַּרְעֹה׃}
}
{
	\eJ{
		And Pharaoh’s officers saw her and praised her to Pharaoh,
		and the woman was taken to Pharaoh’s house.
	}
}
{
	Abe's plan works.
	
	The king sees his hot wife and takes her into his place.
}

\Verse{16}
{
	\hJ{וּלְאַבְרָם הֵיטִיב בַּעֲבוּרָהּ וַיְהִי לוֹ צֹאן וּבָקָר וַחֲמֹרִים וַעֲבָדִים וּשְׁפָחֹת וַאֲתֹנֹת וּגְמַלִּים׃}
}
{
	\eJ{
		And he was good to Abram on her account, and he had a
		flock and oxen and he-asses and servants and maids and
		she-asses and camels.
	}
}
{
	He's really nice to Abe because of her.
	
	He gives him lots of stuff.

	Servants and maids.

	He-asses and She-asses.

	Also animals.
}

\Verse{17}
{
	\hJ{וַיְנַגַּע יְהֹוָה אֶת פַּרְעֹה נְגָעִים גְּדֹלִים וְאֶת בֵּיתוֹ עַל דְּבַר שָׂרַי אֵשֶׁת אַבְרָם׃}
}
{
	\eJ{
		And YHWH plagued Pharaoh and his house, big plagues, over
		the matter of Sarai, Abram’s wife.
	}
}
{
	YHWH, seemingly unaware of the finer points of moral philosophy, gives Pharaoh some plagues for being with Abe's wife.

	Unacceptable, Pharaoh.

	She was his wife.
}

\Verse{18}
{
	\hJ{וַיִּקְרָא פַרְעֹה לְאַבְרָם וַיֹּאמֶר מַה זֹּאת עָשִׂיתָ לִּי לָמָּה לֹא הִגַּדְתָּ לִּי כִּי אִשְׁתְּךָ הִוא׃}
}
{
	\eJ{
		And Pharaoh called Abram and said, “What is this you’ve
		done to me? Why didn’t you tell me that she was your wife?
	}
}
{
	So Pharaoh's like:

	\begin{center}
	\egyptNew{
		𓀠%
		\tikz[baseline=-0.30em]{
			\node[inner sep=0pt] (balls) {\egyptNew{𓨊}};
			\node[
				inner sep=0pt,
				anchor=south,
				yshift=0.10em
			] at (balls.north) {\egyptNew{𓂾}};
		}%
		𓀟𓂿𓀐%
		\tikz[baseline=0.40em]{
			\node[inner sep=0pt] (guy) {\egyptNew{𓀿}};
			\node[
				inner sep=0pt,
				anchor=south,
				yshift=0.10em
			] at (guy.north) {\egyptNew{𓉸}};
		}%
	}\fA{
		Lit. ``Abe! You're breaking my balls man,
		cut it out, you're killin' me.''
	}%
	\aB{\egyptNew{\hspace{-0.1em}𓆤}}%
	\fB{Optional Midrash: I picture him covered in bees here.}
	\end{center}
}

\Verse{19}
{
	\hJ{לָמָה אָמַרְתָּ אֲחֹתִי הִוא וָאֶקַּח אֹתָהּ לִי לְאִשָּׁה וְעַתָּה הִנֵּה אִשְׁתְּךָ קַח וָלֵךְ׃}
}
{
	\eJ{
		Why did you say, ‘She’s my sister,’ so I took her for a
		wife for myself! And now, here’s your wife. Take her and
		go.”
	}
}
{
	I never would have... take your wife and go.

}

\Verse{20}
{
	\hJ{וַיְצַו עָלָיו פַּרְעֹה אֲנָשִׁים וַיְשַׁלְּחוּ אֹתוֹ וְאֶת אִשְׁתּוֹ וְאֶת כּל אֲשֶׁר לוֹ׃}
}
{
	\eJ{
		And Pharaoh commanded people over him, and they let him
		and his wife and all he had go.
	}
}
{
	And they do.
}
